\section{FORMAL RESTATEMENT}
\textbf{Erd\H{o}s \#125; reconstruction, 2026-09-23.}
Let $A$ and $B$ be the nonnegative integers whose base-$3$ and base-$4$
digits, respectively, belong to $\{0,1\}$. Does $A+B$ have positive
lower asymptotic density? We prove that its lower density is zero.
This does not assert that its upper density is zero.

\section{QUICK LITERATURE AND CONTEXT CHECK}
The current problem page records a negative answer and a Lean proof.
The linked formal development uses aligned powers of $3$ and $4$ and a
contraction of counting densities. We reconstruct that argument in
ordinary mathematical notation. This is an exposition of the known
resolution, with no claim of novelty or of having run the Lean checker.

\section{WORK}
Put $S=A+B$, $S(X)=|S\cap[0,X)|$, and $d(X)=S(X)/X$ for integral $X>0$.
Splitting the digit strings after $k$ and $m$ positions gives
\[
 a=3^k a'+a_0,\quad b=4^m b'+b_0,
 \qquad a'\in A,\ b'\in B,
\]
with $0\le a_0\le(3^k-1)/2$ and $0\le b_0\le(4^m-1)/3$.
Let $q=\min(3^k,4^m)$ and $r=\max(3^k,4^m)$. Consequently every
$x=a+b<qM$ can be written
\[
 x=qy+c,\qquad y=a'+b'\in S,\quad 0\le y<M,
\]
where
\[
 0\le c\le C:=(3^k-1)/2+(4^m-1)/3+(r-q)M.       \tag{1}
\]
Indeed the excess term is $(r-q)a'$ or $(r-q)b'$, and either quotient
is at most $y<M$. Counting possible pairs $(y,c)$, without requiring
unique representations, yields
\[
 S(qM)\le S(M)(C+1),\qquad
 d(qM)\le d(M)\frac{C+1}{q}.                    \tag{2}
\]

For any fixed $M$, there are arbitrarily large $k,m$ with $r/q\to1$.
Here is the precise elementary input. The number
$\theta=\log4/\log3$ is irrational: a rational equality would imply
$4^u=3^v$ for positive integers $u,v$. Dirichlet approximation supplies
unbounded $m$ and integers $k$ with $|m\theta-k|\to0$; the denominators
must be unbounded since no fixed nonzero $m$ gives equality. Exponentiating
gives $4^m/3^k\to1$, with both powers tending to infinity.
Along these pairs, (1) gives
\[
 \frac{C+1}{q}\longrightarrow\frac12+\frac13=\frac56.
\]
If $\delta=\liminf_{X\to\infty}d(X)$, (2), with $M$ fixed, therefore implies
\[
 \delta\le\frac56d(M)\quad\hbox{for every positive integer }M.
\]
Choose a sequence $M\to\infty$ on which $d(M)\to\delta$. Then
$\delta\le5\delta/6$, so $\delta=0$.

There is also an explicit iterative formulation. Starting with $M_0=1$,
at stage $j$ select aligned powers with $q_j\ge100$ and
$(r_j-q_j)/q_j\le1/(100M_j)$. Formula (1) gives
\[
 (C_j+1)/q_j\le (5/6)(1+1/(100M_j))+1/100+1/100<9/10.
\]
Set $M_{j+1}=q_jM_j$. Then $M_j\to\infty$ and
$d(M_j)\le(9/10)^j$. Thus the proof produces a subsequence of explicitly
searchable scales at which the density tends to zero.\hfill$\square$

\section{VERIFICATION}
All counting is over $[0,X)$, including $0$; changing finitely many
terms has no effect on density. The upper bound for $c$ is deliberately
allowed to overcount. Alignment must become more accurate as $M$ grows;
a single approximate equality $3^k\approx4^m$ is insufficient for all
iterations. The batch script checks the digit decomposition and (2)
for several small scales.

\section{FINAL STATUS}
\textbf{SOLVED (known negative answer reconstructed).}
The lower density is zero. No conclusion about the existence or value of
the upper density is included. The only general input is elementary
Dirichlet approximation, with its application explained above.

\section{SOURCES}
\url{https://www.erdosproblems.com/125}.
The linked formal-conjectures development, specifically its digit-splitting,
counting, and aligned-scale lemmas:
\url{https://github.com/mo271/formal-conjectures/blob/c27415379b5dbe34105d1fdd707994540c4c6fc7/FormalConjectures/ErdosProblems/125.lean}.
Checked 2026-09-23. Completion measures the lower-density question.

\section{COMPLETION ESTIMATE}
\textbf{COMPLETION: 100\%}
